\documentclass[11pt,a4paper]{article}
\usepackage[margin=24mm]{geometry}
\usepackage{amsmath,amssymb}
\input{glyphtounicode}
\pdfgentounicode=1
\usepackage[hidelinks]{hyperref}
\hypersetup{hypertexnames=false}
\setlength{\parindent}{0pt}
\setlength{\parskip}{6pt}
\begin{document}
\begin{center}
{\Large M2: boundary speed and the first restriction}\\[4pt]
15 October 2026 --- conditional analytical audit, no selected replacement law
\end{center}

\textbf{Status.} The shared-meter M2 direction is user-selected at CP-D2.
This note checks its boundary inputs and activation requirements before
qualifying an exact update. It does not implement a simulator, run traffic,
calibrate parameters, or transfer M1's proofs. Accepted predecessor and M1
source/PDF artifacts remain unchanged.

\section*{1. A flow value does not determine actual boundary speed}

Retain free-speed parameter $v>0$, backward-wave speed $w>0$ and jam
density $k_j>0$. The triangular fundamental diagram is
\begin{equation}
 Q_v(k)=\min\{vk,w(k_j-k)\},\qquad
 k_c(v)=\frac{wk_j}{v+w},\qquad C(v)=vk_c(v).
 \tag{1}
\end{equation}
For any $0<f<C(v)$, it has two density states and associated speeds:
\begin{align}
 k_{\rm f}(f)&=f/v,& u_{\rm f}(f)&=v,\nonumber\\
 k_{\rm c}(f)&=k_j-f/w,&
 u_{\rm c}(f)&=\frac{wf}{wk_j-f}<v.
 \tag{2}
\end{align}
The same subcapacity flow may therefore be light fast traffic or dense
slow traffic. The branch must have a traffic interpretation. Applying
$u_{\rm c}$ to every positive small discharge manufactures a slow state
without establishing that traffic is congested. Empty-flow vehicle speed
is not assigned by this argument.

The inverse service-equivalent speed
\begin{equation}
 u_R=\frac{wR}{wk_j-R},\qquad C(u_R)=R,
 \quad 0<R<C(v)
 \tag{3}
\end{equation}
is likewise not actual crossing speed. A baseline service ceiling below
$C(v)$ does not mean that all undisturbed vehicles travel at $u_R$.
Boundary insertion kinematics require an explicit local regime/entry law.

\textbf{Source interpretation.}
\href{https://arxiv.org/pdf/1405.7080}{Jin's continuous-LTM manuscript}
provides the FD and demand/service interfaces. The
\href{https://www.researchgate.net/publication/251530773_Capacity_Drops_at_Merges_an_endogenous_model}{2011 conference source, section 2},
keeps the main road sufficiently loaded for congestion; its section 2.2
uses main-road speed for an uncongested entering stream. The
\href{https://repository.tudelft.nl/file/File_c424c1b3-7374-436b-a6a7-abdf69063aa6}{2016 manuscript, section 3.2},
relates insertion flow/speed under a congested entering-stream assumption.
Neither relation is an already qualified empty-start M2 update.

\newpage
\section*{2. An undisturbed solution for a speed-recovery-only completion}

Retain the accepted finite-cohort H1 arrivals, $r=e_0+a_0$, continuing
symmetry, initially empty terminal and known zero recovery history.
For fixed advice $v\in[v_-,v_+]$, the minimum-travel eligible arrivals obey
\begin{equation}
 \dot P_A(t)+\dot P_E(t)=vr
 \quad\text{for }t_0=(L-\ell)/v<t<t_*=L/v,
 \qquad 0\text{ otherwise}.
 \tag{4}
\end{equation}
Their composition changes through H1, but their sum does not. The inherited
ordinary budget satisfies
\begin{equation}
 v_+r<R\leq C(v_-),\qquad v_+b_0\leq R_B\leq C(v_-).
 \tag{5}
\end{equation}

Consider the \emph{particular class of completions} with these properties:
(i) zero disturbance gives shared service $R$;
(ii) an unqueued free approach with no separate maneuver burden crosses
at $v$ (or an unqueued entry matches the free continuing stream);
(iii) such nonrestrictive passages generate no recovery loss, and zero
history remains zero in their presence;
(iv) service is conservative and work conserving with the retained priority.
No stochastic disturbance, compulsory-maneuver slowing or new initial state
is included in these conditions.

Then the following trajectory is self-consistent:
\begin{equation}
 Y_A=P_A,\quad Y_E=P_E,\quad Y_B=P_B,\quad
 q_A=q_E=q_B=0,\quad \mu=R,\quad Z=0.
 \tag{6}
\end{equation}
$Z$ denotes the entire recovery state/history, not necessarily a scalar.
Indeed, $\dot P_A+\dot P_E=vr<R$ makes all arrivals pass even under strict
A priority. Each rate is below physical sending capacity. Thus no queue
forms, the speed interpretation remains free, and no new loss is generated.
If a completion with properties (i)--(iv) has a \emph{unique} causal solution,
(6) is that solution. This note does not prove uniqueness for an unspecified
M2 law or rule out other M2 mechanisms. Multiple-solution branch selection
would itself require qualification.

The local complete objective in this completion is
\begin{equation}
 J(v)=\frac{M(L-\ell/2)}{v},\qquad
 J'(v)=-\frac{M(L-\ell/2)}{v^2}<0.
 \tag{7}
\end{equation}
Hence the fastest advice is common and lane composition cannot alter the
objective for fixed aggregate footprint/mass. This is an analytical result
for the stated completion, not an empirical gradient result or a theorem
about every state-dependent shared meter. H1 can still have active early
receiving limitation; consequential terminal waiting is missing.

\newpage
\section*{3. Ordinary overload can be admitted without abandoning H1}

The alternative congestion-origin route must explicitly amend the first
inequality in (5). It need not violate H1's free-flow branch or its terminal
storage protection. To show a nonempty compatible family, choose
\begin{equation}
 0<r<k_c(v_+),\qquad
 v_-r<R<\min\{v_+r,C(v_-)\}.
 \tag{8}
\end{equation}
The interval for $R$ is nonempty: $v_-r<C(v_-)$ follows from
$r<k_c(v_+)\leq k_c(v_-)$, and $v_-r<v_+r$. It permits an ordinary
overload for $v>R/r$ while keeping $R\leq C(v_-)$.

For all admitted advices, H1 still has
\begin{equation}
 e+a=r<k_c(v),\qquad
 \dot F_E,\dot F_A\leq vr<C(v).
 \tag{9}
\end{equation}
Choose $a_0=b_0>0$ sufficiently small, $e_0=r-a_0>0$, and retain
an independent $R_B$ with $v_+b_0\leq R_B\leq C(v_-)$. With total
mass $M=\ell(r+a_0)$ choose finite terminal length $b>M/k_j$.
The earlier vacancy proof then still gives positive receiving at the H1
handoff for any conservative admissible departures:
\begin{equation}
 V_j(t)>0\quad(0<t\leq b/w),\qquad
 V_j(t)\geq k_jb-M>0\quad(t>b/w).
 \tag{10}
\end{equation}
The first part uses $F_j(t)\leq C(v)t<k_jwt$. No spatial H1 state or
upstream spillback feedback is inferred from this sufficient domain.

If a prospective meter has $0\leq\mu\leq R$ and $vr>R$, its combined
eligible queue during the arrival interval must satisfy
\begin{equation}
 Q(t)=P_A(t)+P_E(t)-Y_A(t)-Y_E(t)
 \geq(vr-R)(t-t_0)>0\quad(t_0<t<t_*).
 \tag{11}
\end{equation}
No initial queue or hidden slow-speed seed was inserted. This establishes
a source of terminal congestion, not which stream queues, an insertion
speed law, an extra capacity loss or M2 well-posedness/clearance. Those remain
law-specific. The no-ordinary-overload isolation claim is changed by (8), so
this route requires explicit G3 approval.

\textbf{The other route.} Keeping (5) requires a separately justified
compulsory-entry/yielding mechanism that can generate a restriction before
ordinary overload. Its exact law is not supplied here. Simply reusing M1's
fixed low insertion prototype would reinstate an assumption rather than
derive actual kinematics.

\textbf{Gate.} CP-M2-D1 selects the physical origin of the first restriction:
compulsory-maneuver slowing under the old load domain, or congestion-origin
slowing under an explicitly amended domain. Neither is adopted in this note.
M2 architecture remains selected, with one aggregate account and sufficient
internal information. No M1 theorem, native result or positive control claim
is transferred. Step 3 and the full M2 law review remain open.
\end{document}
